
Bidirectional alignment is operationalized through two observable behavioral signals in multi-turn conversations: user learning rate (reformulation slope) and AI responsiveness (query-response diversity correlation).

\subsubsection{Bidirectional Alignment Framework}

Bidirectional alignment emerges when both participants in a conversation adapt to each other. In human-AI interactions:

\begin{enumerate}
\item \textbf{User learning}: Users discover what queries the AI handles well and refine their communication strategies. Observable signal: reformulation rate decreases over turns.
\end{enumerate}

\begin{enumerate}
\item \textbf{AI responsiveness}: The AI system (if trained with responsiveness mechanisms) adjusts output characteristics based on user query patterns. Observable signal: response diversity correlates with query diversity.
\end{enumerate}

\begin{enumerate}
\item \textbf{Coupling}: When both behaviors co-occur, bidirectional alignment is present. Coupling strength---measured as correlation between user learning rate and AI responsiveness---quantifies alignment quality.
\end{enumerate}

This framework differs from unidirectional alignment (RLHF optimization) by measuring conversation-level dynamics rather than static policy quality.

\subsubsection{Reformulation Detection and Learning Curves}

\textbf{Reformulation definition}: Reformulation is the act of rephrasing a previous query while maintaining semantic intent. A user reformulates when (1) consecutive queries address the same underlying information need, and (2) the phrasing differs syntactically or lexically.

\textbf{Detection algorithm}: For each conversation, consecutive query pairs $(q_i, q_{i+1})$ are extracted and classified as reformulation or novel query using a hybrid semantic-syntactic approach:

\begin{itemize}
\item \textbf{Semantic similarity}: SBERT cosine similarity between query embeddings \citep{reimers2019sentence}. High similarity ($\text{sim} > 0.7$) indicates queries address the same topic.
\item \textbf{Syntactic distance}: Normalized Levenshtein edit distance. High distance ($\text{dist} > 0.3$) indicates different surface forms despite semantic overlap.
\item \textbf{Reformulation criterion}: $(q_i, q_{i+1})$ is a reformulation if $\text{sim}(q_i, q_{i+1}) > 0.7$ AND $\text{edit\_dist}(q_i, q_{i+1}) > 0.3$.
\end{itemize}

This hybrid approach balances precision and recall. Semantic similarity alone produces false positives on topic shifts with overlapping vocabulary. Edit distance alone produces false negatives on paraphrases. Combining both filters out topic shifts while retaining genuine reformulations.

\textbf{Reformulation slope computation}: For each conversation $c$ with $T$ turns:

\begin{enumerate}
\item \textbf{Per-turn reformulation indicator}: $r_t \in \{0, 1\}$ indicating whether query at turn $t$ is a reformulation of query at turn $t-1$.
\item \textbf{Reformulation slope}: Linear regression $r_t \sim \beta_0 + \beta_1 \cdot t$. Extract coefficient $\beta_1$, representing rate of change in reformulation frequency over turns.
\item \textbf{Interpretation}: Negative slope ($\beta_1 < 0$) indicates learning---users reformulate less as they discover effective query strategies. Flat slope ($\beta_1 \approx 0$) indicates no learning. Positive slope ($\beta_1 > 0$) may indicate engagement decay or increasing frustration.
\end{enumerate}

Conversations are filtered to $T \geq 5$ turns, as slope estimation with fewer data points is unreliable. Conversations with insufficient reformulation variation are assigned slope = 0.

\subsubsection{Diversity Metrics and Responsiveness}

\textbf{Diversity operationalization}: Diversity is measured through lexical variation using the Distinct-1 metric \citep{li2016diversity}:

$$\text{Distinct-1}(S) = \frac{|\text{unique unigrams in } S|}{|\text{total unigrams in } S|}$$

This captures vocabulary richness. Distinct-1 is computed separately for user queries and AI responses within each conversation.

\begin{itemize}
\item \textbf{Query diversity}: Distinct-1 across all user queries in conversation $c$.
\item \textbf{Response diversity}: Distinct-1 across all AI responses in conversation $c$.
\end{itemize}

\textbf{Responsiveness measurement}: AI responsiveness is operationalized as the correlation between query diversity and response diversity:

$$\text{Responsiveness}(c) = \text{Pearson}(D_{\text{query}}, D_{\text{response}})$$

Positive correlation ($r > 0$) indicates the AI adjusts response vocabulary based on query vocabulary. This correlation is computed across conversations, not within conversations.

\textbf{Methodological rationale}: Distinct-1 was chosen over semantic diversity metrics for interpretability and computational efficiency. Lexical diversity is transparent, while semantic diversity requires embedding model choices that introduce additional assumptions. Future work may reveal semantic diversity metrics yield stronger correlations, but Distinct-1 provides a conservative baseline.

\subsubsection{Statistical Testing}

\textbf{Hypothesis 1 (User Learning)}:
\begin{itemize}
\item Null hypothesis ($H_0$): Mean reformulation slope across conversations is zero.
\item Alternative hypothesis ($H_1$): Mean reformulation slope is negative.
\item Test: One-sample t-test on slope distribution. Reject $H_0$ if $p < 0.05$ and mean slope significantly negative.
\item Effect size: Cohen's d.
\end{itemize}

\textbf{Hypothesis 2 (AI Responsiveness)}:
\begin{itemize}
\item Null hypothesis ($H_0$): No correlation between query diversity and response diversity ($r = 0$).
\item Alternative hypothesis ($H_1$): Positive correlation ($r > 0.35$, medium effect).
\item Test: Pearson correlation with significance test. Reject $H_0$ if $p < 0.05$ and confidence interval excludes zero.
\item Target threshold: $r > 0.4$ was preregistered as success criterion, but $r > 0.35$ is accepted as weak evidence given lexical diversity may underestimate semantic responsiveness.
\end{itemize}

\subsubsection{Dataset and Filtering}

The HH-RLHF dataset \citep{bai2022constitutional} contains 161,000 conversations between users and an RLHF-trained AI assistant. Each conversation includes multi-turn dialogue, human helpfulness ratings (0-1 scale), and conversation metadata.

\textbf{Filtering criteria}:
\begin{enumerate}
\item Minimum turn count: Conversations with $\geq 5$ turns (required for slope estimation).
\item Complete metadata: Conversations with valid helpfulness ratings and turn structure.
\item Reformulation coverage: Conversations with at least one detected reformulation (to avoid null slopes from all-novel-query conversations).
\end{enumerate}

After filtering, 88 conversations were analyzed for Hypothesis 1 (reformulation slope) and 169,352 conversations for Hypothesis 2 (diversity correlation). The smaller sample for Hypothesis 1 reflects stricter reformulation detection criteria.

\subsubsection{Limitations and Assumptions}

\textbf{Assumption 1 (Learning vs Disengagement)}: Negative reformulation slope may reflect learning OR engagement decay (users stop reformulating because they give up, not because they learned). This is mitigated by stratifying on conversation outcome in future work---learning should predict success, disengagement should not.

\textbf{Assumption 2 (Reformulation Detection Accuracy)}: The SBERT + edit distance heuristic has not been validated against human annotations. False positives and false negatives may bias slope estimates. Conservative thresholds ($\text{sim} > 0.7, \text{dist} > 0.3$) prioritize precision over recall.

\textbf{Assumption 3 (Distinct-1 as Diversity)}: Lexical diversity captures vocabulary variation but not semantic diversity. Two responses with different words but similar meanings yield high Distinct-1 but low semantic diversity. This may underestimate true AI responsiveness.

\textbf{Assumption 4 (Population-Level Correlation)}: Responsiveness is computed as correlation across conversations, not within conversations. This tests whether query diversity predicts response diversity in aggregate but does not measure per-conversation responsiveness.
